\section{Name the fields before interpreting their bytes}
The telemetry schema is local to this exercise. Its outer value is a map with unsigned
integer keys. Three fields are required, and each has a distinct interpretation.

\tagpdfsetup{table/header-rows={1}}
\begin{center}
  \begin{tabular}{rll}
    Key  & Meaning             & Accepted value                                          \\
    \hline
    0    & Acquisition counter & Unsigned integer, all sixty-four bits                   \\
    1    & Device identifier   & Byte string, zero through thirty-two bytes              \\
    2    & Readings            & Array of zero through sixty-four signed 16-bit integers \\
    3..7 & Optional metadata   & A value within the bounded generic subset               \\
  \end{tabular}
\end{center}

The counter is an unsigned acquisition value, not a calendar date. This chapter does
not infer units or convert it to a wall clock. The identifier is binary: a zero byte
inside it does not end the field. A reading is a signed integer code; the schema does
not attach a temperature scale or perform a floating-point conversion.

Accept each map key once, in any order. The outer map therefore has from three through
eight pairs and must contain keys zero, one and two. Optional values must be structurally
valid within this exercise's subset, but they are not copied into the flat record.
The deterministic writer emits only the three required fields. This is a deliberate
projection of the application record, so optional metadata is not promised a serialized
round trip.
